% Conclusion and outlook section for the FINAL report.
% % Conclusion and outlook section for the FINAL report.
% % Conclusion and outlook section for the FINAL report.
% % Conclusion and outlook section for the FINAL report.
% \input{outlook} inside the final report; citations live in refs_outlook.bib
% (merge into the main .bib). Uses \vk{} and \code{} from the report preamble.

\section{Outlook: a Sanskrit hub for other Indian languages}\label{sec:outlook}

\paragraph{The idea.} Most Indian languages share a large Sanskrit-derived
lexical layer, and the region shares many structural features. Indo-Aryan
languages inherited Sanskrit vocabulary through the Prakrits
(\emph{tadbhava}) and re-borrowed it in learned form (\emph{tatsama}).
Dravidian, Munda and other families borrowed it too; the
\emph{tatsama}/\emph{tadbhava} classification is applied to Sanskrit
loanwords across all these families. Bengali, for instance, is estimated at % TODO source: figure from Wikipedia "Tatsama"; replace with a primary reference or drop
about 40\% \emph{tatsama} vocabulary. Structurally, South Asia is a classic
linguistic area \citep{emeneau1956india}: verb-final order, case relations
close to Pāṇini's \emph{kāraka}s, and shared phonological features. This
makes it attractive to treat Sanskrit, or a representation derived from it,
as a \emph{meta-language} through which Indian languages can be related to
one another.

A note on terminology: for Dravidian languages Sanskrit is a source of
borrowing (a superstratum or adstratum), not a substratum. A substratum is
the language of an earlier population that influences an incoming one.
Indeed, Dravidian is often proposed as a \emph{substratum} of Indo-Aryan,
for example as the source of its retroflex consonants
\citep{emeneau1956india}. The accurate phrase is a \emph{shared Sanskrit
lexical stratum}: inherited in Indo-Aryan, borrowed elsewhere.

\paragraph{Prior work.} The idea has a long history in Indian NLP. The
Computational Pāṇinian Grammar framework \citep{bharati1995nlp} uses
\emph{kāraka} relations as a near-interlingual layer. It underlies the
Anusāraka language accessors and the Sampark consortium's system for
translation between Indian languages, which covered 18 language pairs
\citep{anthes2010automated}. Several systems use Sanskrit itself as an
explicit interlingua, mostly in rule-based frameworks. On the statistical
side, \citet{kunchukuttan2017utilizing} showed that subword-level pivot
translation through a \emph{related} pivot language is highly competitive
with direct translation, and that combining several related pivots can
rival a direct model. On the lexical side, IndoWordNet links the wordnets
of 18 Indian languages through shared synset identifiers, built by
expansion from the Hindi WordNet as pivot \citep{bhattacharyya2010indowordnet}.

\paragraph{What this project adds: a semantic hub, not a symbolic
interlingua.} Translating \emph{through} Sanskrit text adds a lossy step
and runs into Sanskrit's small corpus. A lighter and testable version of the
idea uses our lemma- and sense-level Sanskrit space as the \emph{hub} of a
star-shaped alignment of embedding spaces. Let $X_0\in\mathbb{R}^{n_0\times
d}$ be the Sanskrit lemma space and $X_\ell$ the space of language
$\ell=1,\dots,L$. A seed dictionary $D_\ell$ of \emph{tatsama} pairs is
almost free to construct: after transliteration into one phonemic alphabet,
a \emph{tatsama} is (nearly) string-identical to its Sanskrit lemma. The
transliteration problem of \S\ref{sec:audit} returns here at the scale of
a dozen scripts. The orthogonal map aligning language $\ell$ to the hub,
\begin{align*}
W_\ell&=\arg\min_{W^\top W=I}\big\|X_\ell[D_\ell]\,W-X_0[D_\ell]\big\|_F\\
&=U_\ell V_\ell^\top,
\end{align*}
has the closed form given by the SVD $U_\ell\Sigma_\ell V_\ell^\top$ of
$X_\ell[D_\ell]^\top X_0[D_\ell]$ \citep{schonemann1966generalized}. This is
the standard tool for bilingual embedding alignment
\citep{mikolov2013exploiting,smith2017offline,conneau2018word}. Because
$W_\ell$ is orthogonal it is an isometry, so each language's internal
cosine geometry is untouched. Language $\ell$ is mapped to language $m$ by
$W_\ell W_m^\top$, so $L$ alignments serve all $L(L-1)$ directed pairs.

\paragraph{Why the hub should be sense-level.} A shared lexical layer helps
only to the extent that shared forms share meanings. Daughter languages
select different senses of the same Sanskrit word. \emph{Śikṣā} is
``education'' in Hindi but ``punishment'' in Telugu. \emph{Upanyāsa} (in
Sanskrit roughly ``statement, introduction'') became ``novel'' in Hindi and
``lecture'' in Telugu. These are cross-lingual versions of the within-lemma
polysemy we measured (0.94 bits/token beyond lemma IDs on the sense-annotated
DCS). The prediction follows directly from RQ1 and RQ3: a sense-level hub
(\vk{lid}, \vk{sense}) should align better than a surface or lemma-string
hub, and the alignment residual of a seed pair,
$\|\mathbf{x}_\ell W_\ell-\mathbf{x}_0\|$, should flag semantically drifted
\emph{tatsama}s, i.e.\ false friends. That gives a data-driven detector of
cross-lingual semantic change. The Amarakośa adds a language-neutral concept
inventory: its vargas and synsets could serve as a Sanskrit-anchored
counterpart to IndoWordNet's Hindi-anchored synsets. This would be
particularly natural for Dravidian languages, whose learned vocabulary
borrowed from Sanskrit directly.

\paragraph{Testable hypotheses.}
\begin{enumerate}
\item Alignment quality to the Sanskrit hub (precision@$k$ of word
  translation) increases with a language's \emph{tatsama} share and with the
  formality of the target-side corpus.
\item A sense-level hub outperforms a lemma-string hub, which outperforms a
  surface hub, on held-out \emph{tatsama} and non-\emph{tatsama} pairs.
\item Pairs with large alignment residuals are enriched for documented false
  friends.
\end{enumerate}

\paragraph{Limits.} The shared layer is concentrated in learned and formal
registers. Colloquial language, the core grammatical vocabulary of Dravidian
languages, and Persian- and English-derived vocabulary fall outside it.
Tamil in particular uses far less \emph{tatsama} vocabulary than Telugu,
Kannada or Malayalam. Structural similarity is largely an areal phenomenon,
not inheritance, so it supports a shared \emph{framework} (as in
Computational Pāṇinian Grammar) rather than translation through Sanskrit
itself. Finally, modern multilingual neural translation already pools
related languages implicitly. The case for a Sanskrit hub is therefore
strongest where data are scarcest: low-resource language pairs,
interpretable lexical alignment, and the study of how a shared lexicon has
diverged in meaning.
 inside the final report; citations live in refs_outlook.bib
% (merge into the main .bib). Uses \vk{} and \code{} from the report preamble.

\section{Outlook: a Sanskrit hub for other Indian languages}\label{sec:outlook}

\paragraph{The idea.} Most Indian languages share a large Sanskrit-derived
lexical layer, and the region shares many structural features. Indo-Aryan
languages inherited Sanskrit vocabulary through the Prakrits
(\emph{tadbhava}) and re-borrowed it in learned form (\emph{tatsama}).
Dravidian, Munda and other families borrowed it too; the
\emph{tatsama}/\emph{tadbhava} classification is applied to Sanskrit
loanwords across all these families. Bengali, for instance, is estimated at % TODO source: figure from Wikipedia "Tatsama"; replace with a primary reference or drop
about 40\% \emph{tatsama} vocabulary. Structurally, South Asia is a classic
linguistic area \citep{emeneau1956india}: verb-final order, case relations
close to Pāṇini's \emph{kāraka}s, and shared phonological features. This
makes it attractive to treat Sanskrit, or a representation derived from it,
as a \emph{meta-language} through which Indian languages can be related to
one another.

A note on terminology: for Dravidian languages Sanskrit is a source of
borrowing (a superstratum or adstratum), not a substratum. A substratum is
the language of an earlier population that influences an incoming one.
Indeed, Dravidian is often proposed as a \emph{substratum} of Indo-Aryan,
for example as the source of its retroflex consonants
\citep{emeneau1956india}. The accurate phrase is a \emph{shared Sanskrit
lexical stratum}: inherited in Indo-Aryan, borrowed elsewhere.

\paragraph{Prior work.} The idea has a long history in Indian NLP. The
Computational Pāṇinian Grammar framework \citep{bharati1995nlp} uses
\emph{kāraka} relations as a near-interlingual layer. It underlies the
Anusāraka language accessors and the Sampark consortium's system for
translation between Indian languages, which covered 18 language pairs
\citep{anthes2010automated}. Several systems use Sanskrit itself as an
explicit interlingua, mostly in rule-based frameworks. On the statistical
side, \citet{kunchukuttan2017utilizing} showed that subword-level pivot
translation through a \emph{related} pivot language is highly competitive
with direct translation, and that combining several related pivots can
rival a direct model. On the lexical side, IndoWordNet links the wordnets
of 18 Indian languages through shared synset identifiers, built by
expansion from the Hindi WordNet as pivot \citep{bhattacharyya2010indowordnet}.

\paragraph{What this project adds: a semantic hub, not a symbolic
interlingua.} Translating \emph{through} Sanskrit text adds a lossy step
and runs into Sanskrit's small corpus. A lighter and testable version of the
idea uses our lemma- and sense-level Sanskrit space as the \emph{hub} of a
star-shaped alignment of embedding spaces. Let $X_0\in\mathbb{R}^{n_0\times
d}$ be the Sanskrit lemma space and $X_\ell$ the space of language
$\ell=1,\dots,L$. A seed dictionary $D_\ell$ of \emph{tatsama} pairs is
almost free to construct: after transliteration into one phonemic alphabet,
a \emph{tatsama} is (nearly) string-identical to its Sanskrit lemma. The
transliteration problem of \S\ref{sec:audit} returns here at the scale of
a dozen scripts. The orthogonal map aligning language $\ell$ to the hub,
\begin{align*}
W_\ell&=\arg\min_{W^\top W=I}\big\|X_\ell[D_\ell]\,W-X_0[D_\ell]\big\|_F\\
&=U_\ell V_\ell^\top,
\end{align*}
has the closed form given by the SVD $U_\ell\Sigma_\ell V_\ell^\top$ of
$X_\ell[D_\ell]^\top X_0[D_\ell]$ \citep{schonemann1966generalized}. This is
the standard tool for bilingual embedding alignment
\citep{mikolov2013exploiting,smith2017offline,conneau2018word}. Because
$W_\ell$ is orthogonal it is an isometry, so each language's internal
cosine geometry is untouched. Language $\ell$ is mapped to language $m$ by
$W_\ell W_m^\top$, so $L$ alignments serve all $L(L-1)$ directed pairs.

\paragraph{Why the hub should be sense-level.} A shared lexical layer helps
only to the extent that shared forms share meanings. Daughter languages
select different senses of the same Sanskrit word. \emph{Śikṣā} is
``education'' in Hindi but ``punishment'' in Telugu. \emph{Upanyāsa} (in
Sanskrit roughly ``statement, introduction'') became ``novel'' in Hindi and
``lecture'' in Telugu. These are cross-lingual versions of the within-lemma
polysemy we measured (0.94 bits/token beyond lemma IDs on the sense-annotated
DCS). The prediction follows directly from RQ1 and RQ3: a sense-level hub
(\vk{lid}, \vk{sense}) should align better than a surface or lemma-string
hub, and the alignment residual of a seed pair,
$\|\mathbf{x}_\ell W_\ell-\mathbf{x}_0\|$, should flag semantically drifted
\emph{tatsama}s, i.e.\ false friends. That gives a data-driven detector of
cross-lingual semantic change. The Amarakośa adds a language-neutral concept
inventory: its vargas and synsets could serve as a Sanskrit-anchored
counterpart to IndoWordNet's Hindi-anchored synsets. This would be
particularly natural for Dravidian languages, whose learned vocabulary
borrowed from Sanskrit directly.

\paragraph{Testable hypotheses.}
\begin{enumerate}
\item Alignment quality to the Sanskrit hub (precision@$k$ of word
  translation) increases with a language's \emph{tatsama} share and with the
  formality of the target-side corpus.
\item A sense-level hub outperforms a lemma-string hub, which outperforms a
  surface hub, on held-out \emph{tatsama} and non-\emph{tatsama} pairs.
\item Pairs with large alignment residuals are enriched for documented false
  friends.
\end{enumerate}

\paragraph{Limits.} The shared layer is concentrated in learned and formal
registers. Colloquial language, the core grammatical vocabulary of Dravidian
languages, and Persian- and English-derived vocabulary fall outside it.
Tamil in particular uses far less \emph{tatsama} vocabulary than Telugu,
Kannada or Malayalam. Structural similarity is largely an areal phenomenon,
not inheritance, so it supports a shared \emph{framework} (as in
Computational Pāṇinian Grammar) rather than translation through Sanskrit
itself. Finally, modern multilingual neural translation already pools
related languages implicitly. The case for a Sanskrit hub is therefore
strongest where data are scarcest: low-resource language pairs,
interpretable lexical alignment, and the study of how a shared lexicon has
diverged in meaning.
 inside the final report; citations live in refs_outlook.bib
% (merge into the main .bib). Uses \vk{} and \code{} from the report preamble.

\section{Outlook: a Sanskrit hub for other Indian languages}\label{sec:outlook}

\paragraph{The idea.} Most Indian languages share a large Sanskrit-derived
lexical layer, and the region shares many structural features. Indo-Aryan
languages inherited Sanskrit vocabulary through the Prakrits
(\emph{tadbhava}) and re-borrowed it in learned form (\emph{tatsama}).
Dravidian, Munda and other families borrowed it too; the
\emph{tatsama}/\emph{tadbhava} classification is applied to Sanskrit
loanwords across all these families. Bengali, for instance, is estimated at % TODO source: figure from Wikipedia "Tatsama"; replace with a primary reference or drop
about 40\% \emph{tatsama} vocabulary. Structurally, South Asia is a classic
linguistic area \citep{emeneau1956india}: verb-final order, case relations
close to Pāṇini's \emph{kāraka}s, and shared phonological features. This
makes it attractive to treat Sanskrit, or a representation derived from it,
as a \emph{meta-language} through which Indian languages can be related to
one another.

A note on terminology: for Dravidian languages Sanskrit is a source of
borrowing (a superstratum or adstratum), not a substratum. A substratum is
the language of an earlier population that influences an incoming one.
Indeed, Dravidian is often proposed as a \emph{substratum} of Indo-Aryan,
for example as the source of its retroflex consonants
\citep{emeneau1956india}. The accurate phrase is a \emph{shared Sanskrit
lexical stratum}: inherited in Indo-Aryan, borrowed elsewhere.

\paragraph{Prior work.} The idea has a long history in Indian NLP. The
Computational Pāṇinian Grammar framework \citep{bharati1995nlp} uses
\emph{kāraka} relations as a near-interlingual layer. It underlies the
Anusāraka language accessors and the Sampark consortium's system for
translation between Indian languages, which covered 18 language pairs
\citep{anthes2010automated}. Several systems use Sanskrit itself as an
explicit interlingua, mostly in rule-based frameworks. On the statistical
side, \citet{kunchukuttan2017utilizing} showed that subword-level pivot
translation through a \emph{related} pivot language is highly competitive
with direct translation, and that combining several related pivots can
rival a direct model. On the lexical side, IndoWordNet links the wordnets
of 18 Indian languages through shared synset identifiers, built by
expansion from the Hindi WordNet as pivot \citep{bhattacharyya2010indowordnet}.

\paragraph{What this project adds: a semantic hub, not a symbolic
interlingua.} Translating \emph{through} Sanskrit text adds a lossy step
and runs into Sanskrit's small corpus. A lighter and testable version of the
idea uses our lemma- and sense-level Sanskrit space as the \emph{hub} of a
star-shaped alignment of embedding spaces. Let $X_0\in\mathbb{R}^{n_0\times
d}$ be the Sanskrit lemma space and $X_\ell$ the space of language
$\ell=1,\dots,L$. A seed dictionary $D_\ell$ of \emph{tatsama} pairs is
almost free to construct: after transliteration into one phonemic alphabet,
a \emph{tatsama} is (nearly) string-identical to its Sanskrit lemma. The
transliteration problem of \S\ref{sec:audit} returns here at the scale of
a dozen scripts. The orthogonal map aligning language $\ell$ to the hub,
\begin{align*}
W_\ell&=\arg\min_{W^\top W=I}\big\|X_\ell[D_\ell]\,W-X_0[D_\ell]\big\|_F\\
&=U_\ell V_\ell^\top,
\end{align*}
has the closed form given by the SVD $U_\ell\Sigma_\ell V_\ell^\top$ of
$X_\ell[D_\ell]^\top X_0[D_\ell]$ \citep{schonemann1966generalized}. This is
the standard tool for bilingual embedding alignment
\citep{mikolov2013exploiting,smith2017offline,conneau2018word}. Because
$W_\ell$ is orthogonal it is an isometry, so each language's internal
cosine geometry is untouched. Language $\ell$ is mapped to language $m$ by
$W_\ell W_m^\top$, so $L$ alignments serve all $L(L-1)$ directed pairs.

\paragraph{Why the hub should be sense-level.} A shared lexical layer helps
only to the extent that shared forms share meanings. Daughter languages
select different senses of the same Sanskrit word. \emph{Śikṣā} is
``education'' in Hindi but ``punishment'' in Telugu. \emph{Upanyāsa} (in
Sanskrit roughly ``statement, introduction'') became ``novel'' in Hindi and
``lecture'' in Telugu. These are cross-lingual versions of the within-lemma
polysemy we measured (0.94 bits/token beyond lemma IDs on the sense-annotated
DCS). The prediction follows directly from RQ1 and RQ3: a sense-level hub
(\vk{lid}, \vk{sense}) should align better than a surface or lemma-string
hub, and the alignment residual of a seed pair,
$\|\mathbf{x}_\ell W_\ell-\mathbf{x}_0\|$, should flag semantically drifted
\emph{tatsama}s, i.e.\ false friends. That gives a data-driven detector of
cross-lingual semantic change. The Amarakośa adds a language-neutral concept
inventory: its vargas and synsets could serve as a Sanskrit-anchored
counterpart to IndoWordNet's Hindi-anchored synsets. This would be
particularly natural for Dravidian languages, whose learned vocabulary
borrowed from Sanskrit directly.

\paragraph{Testable hypotheses.}
\begin{enumerate}
\item Alignment quality to the Sanskrit hub (precision@$k$ of word
  translation) increases with a language's \emph{tatsama} share and with the
  formality of the target-side corpus.
\item A sense-level hub outperforms a lemma-string hub, which outperforms a
  surface hub, on held-out \emph{tatsama} and non-\emph{tatsama} pairs.
\item Pairs with large alignment residuals are enriched for documented false
  friends.
\end{enumerate}

\paragraph{Limits.} The shared layer is concentrated in learned and formal
registers. Colloquial language, the core grammatical vocabulary of Dravidian
languages, and Persian- and English-derived vocabulary fall outside it.
Tamil in particular uses far less \emph{tatsama} vocabulary than Telugu,
Kannada or Malayalam. Structural similarity is largely an areal phenomenon,
not inheritance, so it supports a shared \emph{framework} (as in
Computational Pāṇinian Grammar) rather than translation through Sanskrit
itself. Finally, modern multilingual neural translation already pools
related languages implicitly. The case for a Sanskrit hub is therefore
strongest where data are scarcest: low-resource language pairs,
interpretable lexical alignment, and the study of how a shared lexicon has
diverged in meaning.
 inside the final report; citations live in refs_outlook.bib
% (merge into the main .bib). Uses \vk{} and \code{} from the report preamble.

\section{Outlook: a Sanskrit hub for other Indian languages}\label{sec:outlook}

\paragraph{The idea.} Most Indian languages share a large Sanskrit-derived
lexical layer, and the region shares many structural features. Indo-Aryan
languages inherited Sanskrit vocabulary through the Prakrits
(\emph{tadbhava}) and re-borrowed it in learned form (\emph{tatsama}).
Dravidian, Munda and other families borrowed it too; the
\emph{tatsama}/\emph{tadbhava} classification is applied to Sanskrit
loanwords across all these families. Bengali, for instance, is estimated at % TODO source: figure from Wikipedia "Tatsama"; replace with a primary reference or drop
about 40\% \emph{tatsama} vocabulary. Structurally, South Asia is a classic
linguistic area \citep{emeneau1956india}: verb-final order, case relations
close to Pāṇini's \emph{kāraka}s, and shared phonological features. This
makes it attractive to treat Sanskrit, or a representation derived from it,
as a \emph{meta-language} through which Indian languages can be related to
one another.

A note on terminology: for Dravidian languages Sanskrit is a source of
borrowing (a superstratum or adstratum), not a substratum. A substratum is
the language of an earlier population that influences an incoming one.
Indeed, Dravidian is often proposed as a \emph{substratum} of Indo-Aryan,
for example as the source of its retroflex consonants
\citep{emeneau1956india}. The accurate phrase is a \emph{shared Sanskrit
lexical stratum}: inherited in Indo-Aryan, borrowed elsewhere.

\paragraph{Prior work.} The idea has a long history in Indian NLP. The
Computational Pāṇinian Grammar framework \citep{bharati1995nlp} uses
\emph{kāraka} relations as a near-interlingual layer. It underlies the
Anusāraka language accessors and the Sampark consortium's system for
translation between Indian languages, which covered 18 language pairs
\citep{anthes2010automated}. Several systems use Sanskrit itself as an
explicit interlingua, mostly in rule-based frameworks. On the statistical
side, \citet{kunchukuttan2017utilizing} showed that subword-level pivot
translation through a \emph{related} pivot language is highly competitive
with direct translation, and that combining several related pivots can
rival a direct model. On the lexical side, IndoWordNet links the wordnets
of 18 Indian languages through shared synset identifiers, built by
expansion from the Hindi WordNet as pivot \citep{bhattacharyya2010indowordnet}.

\paragraph{What this project adds: a semantic hub, not a symbolic
interlingua.} Translating \emph{through} Sanskrit text adds a lossy step
and runs into Sanskrit's small corpus. A lighter and testable version of the
idea uses our lemma- and sense-level Sanskrit space as the \emph{hub} of a
star-shaped alignment of embedding spaces. Let $X_0\in\mathbb{R}^{n_0\times
d}$ be the Sanskrit lemma space and $X_\ell$ the space of language
$\ell=1,\dots,L$. A seed dictionary $D_\ell$ of \emph{tatsama} pairs is
almost free to construct: after transliteration into one phonemic alphabet,
a \emph{tatsama} is (nearly) string-identical to its Sanskrit lemma. The
transliteration problem of \S\ref{sec:audit} returns here at the scale of
a dozen scripts. The orthogonal map aligning language $\ell$ to the hub,
\begin{align*}
W_\ell&=\arg\min_{W^\top W=I}\big\|X_\ell[D_\ell]\,W-X_0[D_\ell]\big\|_F\\
&=U_\ell V_\ell^\top,
\end{align*}
has the closed form given by the SVD $U_\ell\Sigma_\ell V_\ell^\top$ of
$X_\ell[D_\ell]^\top X_0[D_\ell]$ \citep{schonemann1966generalized}. This is
the standard tool for bilingual embedding alignment
\citep{mikolov2013exploiting,smith2017offline,conneau2018word}. Because
$W_\ell$ is orthogonal it is an isometry, so each language's internal
cosine geometry is untouched. Language $\ell$ is mapped to language $m$ by
$W_\ell W_m^\top$, so $L$ alignments serve all $L(L-1)$ directed pairs.

\paragraph{Why the hub should be sense-level.} A shared lexical layer helps
only to the extent that shared forms share meanings. Daughter languages
select different senses of the same Sanskrit word. \emph{Śikṣā} is
``education'' in Hindi but ``punishment'' in Telugu. \emph{Upanyāsa} (in
Sanskrit roughly ``statement, introduction'') became ``novel'' in Hindi and
``lecture'' in Telugu. These are cross-lingual versions of the within-lemma
polysemy we measured (0.94 bits/token beyond lemma IDs on the sense-annotated
DCS). The prediction follows directly from RQ1 and RQ3: a sense-level hub
(\vk{lid}, \vk{sense}) should align better than a surface or lemma-string
hub, and the alignment residual of a seed pair,
$\|\mathbf{x}_\ell W_\ell-\mathbf{x}_0\|$, should flag semantically drifted
\emph{tatsama}s, i.e.\ false friends. That gives a data-driven detector of
cross-lingual semantic change. The Amarakośa adds a language-neutral concept
inventory: its vargas and synsets could serve as a Sanskrit-anchored
counterpart to IndoWordNet's Hindi-anchored synsets. This would be
particularly natural for Dravidian languages, whose learned vocabulary
borrowed from Sanskrit directly.

\paragraph{Testable hypotheses.}
\begin{enumerate}
\item Alignment quality to the Sanskrit hub (precision@$k$ of word
  translation) increases with a language's \emph{tatsama} share and with the
  formality of the target-side corpus.
\item A sense-level hub outperforms a lemma-string hub, which outperforms a
  surface hub, on held-out \emph{tatsama} and non-\emph{tatsama} pairs.
\item Pairs with large alignment residuals are enriched for documented false
  friends.
\end{enumerate}

\paragraph{Limits.} The shared layer is concentrated in learned and formal
registers. Colloquial language, the core grammatical vocabulary of Dravidian
languages, and Persian- and English-derived vocabulary fall outside it.
Tamil in particular uses far less \emph{tatsama} vocabulary than Telugu,
Kannada or Malayalam. Structural similarity is largely an areal phenomenon,
not inheritance, so it supports a shared \emph{framework} (as in
Computational Pāṇinian Grammar) rather than translation through Sanskrit
itself. Finally, modern multilingual neural translation already pools
related languages implicitly. The case for a Sanskrit hub is therefore
strongest where data are scarcest: low-resource language pairs,
interpretable lexical alignment, and the study of how a shared lexicon has
diverged in meaning.
